% !TITLE 长相思·其二
% !AUTHOR 李白
% !DYNASTY 唐
% !GENRE 乐府诗
% !ORDER 15

\begin{poem}
日色欲尽花含烟\textsuperscript{1}，月明如素\textsuperscript{2}愁不眠。
赵瑟\textsuperscript{3}初停凤凰柱，蜀琴\textsuperscript{4}欲奏鸳鸯弦。
此曲有意无人传，愿随春风寄燕然\textsuperscript{5}。
忆君迢迢隔青天。
昔时横波目\textsuperscript{6}，今作流泪泉。
不信妾肠断，归来看取明镜前。
\end{poem}

\begin{annotations}
\notetext{1}{花含烟：黄昏水汽中花朵若隐若现，像含着一层轻烟。}
\notetext{2}{月明如素：月光明亮如白色的绢。素，白色生绢。}
\notetext{3}{赵瑟：赵地（今河北一带）出产的瑟，以音质优美著称。凤凰柱，瑟上雕有凤凰形状的弦柱。}
\notetext{4}{蜀琴：蜀地（今四川）出产的琴。鸳鸯弦，琴弦成双成对，反衬弹琴者的孤独。}
\notetext{5}{燕然：燕然山，在今蒙古国境内。代指遥远的边疆，心上人戍守的地方。}
\notetext{6}{横波目：眼神像水波一样流转的眼睛，形容美目的灵动。如今这双眼睛变成了流不干的泪泉。}
\end{annotations}

\begin{appreciation}
这是《长相思》的第二首，从女子的视角写相思。如果说第一首是男子对远方美人的遥望，这一首就是女子对远方男子的倾诉。

日色欲尽花含烟，月明如素愁不眠——太阳要落了，花儿笼罩在烟霭中；月亮明亮如白绢，她却愁得无法入睡。花含烟是一个极美的意象：黄昏时分，水汽氤氲，花朵若隐若现，像含着一层轻烟。这种朦胧的美，和女子内心的忧愁形成了微妙的对应。

赵瑟初停凤凰柱，蜀琴欲奏鸳鸯弦——她刚停下了赵瑟，又想弹奏蜀琴。瑟和琴都是古代的弦乐器。凤凰柱、鸳鸯弦，这些名字本身就带有成双成对的意味，更衬出她独自一人的孤寂。此曲有意无人传，愿随春风寄燕然——这首曲子有情意，却没有人能传达；只愿随着春风，寄到燕然山。燕然山在今蒙古国境内，是边塞的代称。她的心上人在边疆。

忆君迢迢隔青天——回忆你啊，隔着迢迢青天。昔时横波目，今作流泪泉——昔日那流转如波的眼睛，如今变成了流泪的泉水。横波目形容眼神流动如水波，是古诗中形容美目的经典词汇。但现在，这双美丽的眼睛只剩下了泪水。不信妾肠断，归来看取明镜前——你若不信我肝肠寸断，就回来照照镜子看看我。镜子不会说谎，镜中的容颜会告诉她思念的深度。

这两首《长相思》，一首从男子角度，一首从女子角度，构成了一个完整的对话。但他们永远不会真正对话——因为他们隔着关山，隔着时间，隔着命运。李白的伟大之处在于，他写的是个人的相思，却写出了人类普遍的孤独。那种无论如何都无法到达对方的绝望，是爱情中最深的痛苦，也是生命中最真的体验。
\end{appreciation}
